\scanpage{61}
\mspage{II. 9}
\diagram{\begin{tikzcd}[ampersand replacement=\&]
\operatorname{Gal}(\unsure{\overline K'|\overline K}/K')\arrow[r]\arrow[d] \& \operatorname{Gal}(\overline k'/k')\\
\operatorname{Gal}(\overline K/K) \&
\end{tikzcd}}

% lines 4-16
ont des images qui sont des sous-groupes ouverts. Donc \deleted{l'hyp.} dire que $A_K$, ou $A_{K'}$, ou $B$ sur $k'$, donnent lieu à des actions ess\textsuperscript{t} \deleted{\unsure{résolubles}} \deleted{abéliennes} \inserted{abéliennes}, revient au même (Ceci est vrai chaque fois qu'on a une extension de type fini\unsure{\ldots}) ! \deleted{\unsure{Donc}} On en conclut p.ex. déjà une réciproque \deleted{\unsure{donc la}} de l'énoncé donné, lorsque $\operatorname{car} K>0$ \inserted{donc $k'$ est un corps fini} : s'il existe $B/k'$ avec $B\otimes_{k'}K'$ isogène à $A\otimes_K K'$, alors l'action de \deleted{$\pi(K/K')$ sur} $\operatorname{Gal}(\overline K/K)$ \deleted{sur} $T_\ell(A(\overline K))$ est \deleted{\unsure{alors}} ess\textsuperscript{t} abélienne : en effet, celle de \deleted{\unsure{Gal}} \deleted{\unsure{$\xi$}} $\operatorname{Gal}(\overline k'/k')$ \inserted{$\simeq\widehat{\mathbb{Z}}$} sur $T_\ell(B(\overline k'))$ l'est évidemment !

% lines 17-18
Dans le cas en car $K=0$ cela nous montre que

% lines 19-29
Pour prouver la dernière assertion de 1.9, on est ramené \deleted{\unsure{à pr.}} au cas où $K$ est déjà un corps de nombres. Dans ce cas, \deleted{d'après le} \deleted{\unsure{théorème}} \deleted{\unsure{critère de}} ce résultat est bien connu. Par \deleted{\unsure{le critère de Néron}} \deleted{\unsure{par}} le critère de Néron-\inserted{Ogg}-\unsure{Safarevitch|Schafarevitch} (\unsure{[\quad \S\ S., th.\ 2]|[\quad \S\ S., T.\ 2]}, cf aussi Exp.\quad) on est ramené à voir que si $R$ est l'anneau des entiers de $K$, alors \inserted{pour tt idéal maximal $\mathfrak m$ de $R$,} l'action du groupe d'inertie $I\subset\pi=\operatorname{Gal}(\overline K/K)$ relatif à $R_{\mathfrak m}$ est quasi-\deleted{unipotente}\inserted{triviale} (de sorte qu'après une extension finie convenable \inserted{$K'$} de $K$, on
\margin{\uline{développer}\ \uline{un peu}}
% slanted
% archive sheet number: 60
